\documentclass[10pt,a4paper]{amsart}
\usepackage[margin=19mm]{geometry}
\usepackage{amsmath,amssymb,mathtools}
\usepackage[scr=rsfs,cal=euler]{mathalfa}
\usepackage{xfrac}
\usepackage[unicode,hidelinks,bookmarks=false]{hyperref}
\newcommand{\ie}{i.e.\ }
\newcommand{\cf}{cf.\ }
\newcommand{\resp}{respectively\ }
\newcommand{\Subsection}{\subsection}
\providecommand{\nobreakdash}{\nobreak\hbox{-}}
\setlength{\emergencystretch}{2em}
\multlinegap0pt
\usepackage{amssymb}
\usepackage{mathtools}
\usepackage{enumerate}
\usepackage[shortlabels]{enumitem}
\usepackage{bbm}
\usepackage{xcolor}
\newtheorem{theorem}{Theorem}
\newtheorem{proposition}{Proposition}
\newcommand{\R}{\mathbb{R}}
\newcommand{\probmsymr}{\mathcal{M}^{sym}(\R)}
\title{Functional inequalities for Boolean entropy}
\author{Guillaume C\'ebron, Kewei Pan}
\date{Source: arXiv:2602.23527v1 (CC BY 4.0)}
\begin{document}
\maketitle

\noindent\emph{Source and licence.} Functional inequalities for Boolean entropy. Version arXiv:2602.23527v1 (CC BY 4.0), distributed by arXiv under the Creative Commons Attribution 4.0 International licence, http://creativecommons.org/licenses/by/4.0/, as recorded in the arXiv licence metadata for this exact version and shown on the abstract page https://arxiv.org/abs/2602.23527v1. Full licence notice: \texttt{licenses/arxiv-2602.23527-CC-BY-4.0.txt}.\par\smallskip

\noindent\textit{Editorial scope.} Counted unit: the article's proposition proposition (source0.tex lines 1404--1413). The article's complete original proof of it is delivered with it (source0.tex lines 1415--1490, one \texttt{proof} environment of 76 lines). No mathematics is written, repaired or re-derived: the statement and the proof are the article's own lines, in the article's own notation, and no retained line is substituted or re-flowed.  Retained uncounted as the source-local context the proof needs: lines 543--545; lines 1013--1024; lines 1312--1315.  Nothing was omitted from any retained span, so the package carries no omission record.  No retained line contains a citation, so the wrapper carries no bibliography and no bibliography entry.  No retained line contains a figure environment, an image-inclusion command or drawing code, so no image is delivered and the package declares no asset dependency.  Editorial formatting only, in addition: an AMS \texttt{amsart} preamble that restates the article's own declarations and macros used by the retained lines, namely the environments \texttt{theorem}, \texttt{proposition} on separate counters, together with the standard packages those lines need.  The retained environments print in this document as Theorem 1, Proposition 1 in order of appearance, whereas the article prints the same items with the numbering of its own full document.  The complete native source of the article as distributed in the arXiv e-print for this exact version is bundled unchanged as \texttt{sources/195408/source0.tex}.
		\begin{equation}\label{eq:fisher_0^1}
			\Psi(X) = \iint \frac{\log x^2-\log y^2}{x^2-y^2}\,d\mu(x)d\mu(y) = \iint \int_0^1 \frac{1}{tx^2 + (1-t)y^2}\,dt\,d\mu(x)d\mu(y),
		\end{equation}

	    \begin{theorem}[Constancy in the Central Limit Theorem]\label{th:constancy}
        Let $\{X_i\}_{i\ge 1}\subset (\mathcal{M},\tau)$ be a sequence of i.i.d.\ (in the Boolean sense) random variables with law $\mu\in\probmsymr $. Define
        \begin{equation*}
            S_n:=\frac{X_1+\cdots+X_n}{\sqrt{n}}.
        \end{equation*}
        Then, for any $n\geq 1$, one has
        \[
        \Gamma^*(S_n)= \Gamma^*(X_1)
        \qquad\text{and}\qquad
        \Psi^*(S_n)= \Psi^*(X_1).
        \]
    \end{theorem}

\begin{proposition}[Berry-Esseen bound]\label{prop:Berryesseen}Let $\mu\in \mathcal{M}^{sym}(\mathbb{R})$ such that $m_1(\mu) = 0$, $m_2(\mu) = 1$,
and $m_4(\mu) < +\infty$.  We define $\mu_n$ to be the dilation of $\mu^{\uplus n}$ by $\frac{1}{\sqrt{n}}$, i.e. $\mu_n(B)=\mu^{\uplus n}(\{x:\frac{1}{\sqrt{n}}x\in B\})$. We have
    $$W_2(\mu_n,\mathbb{\mathrm{b}})\leq D^*(\mu_n|\mathbb{\mathrm{b}})=\frac{1}{\sqrt{n}}\sqrt{m_4(\mu)-1}.$$
\end{proposition}

\begin{proposition}
Let $\mu\in \mathcal{M}(\mathbb{R})$ such that $m_1(\mu)=0$, $m_2(\mu)=1$ and $m_4(\mu)<+\infty$.  
Let $\mu_n$ be the dilation of $\mu^{\uplus n}$ by $1/\sqrt{n}$, i, i.e. $\mu_n(B)=\mu^{\uplus n}(\{x:\frac{1}{\sqrt{n}}x\in B\})$.
Then
\[
\Psi(\mu_n|\mathbb{\mathrm b})
\le \Psi^*(\mu_n)\,D^*(\mu_n|\mathbb{\mathrm b})
=\frac{1}{\sqrt n}\,m_{-2}(\mu)\sqrt{m_4(\mu)-1}.
\]
\end{proposition}

\begin{proof}
We first prove the general bound
\[
|\Psi(\mu)-1|
\le \Psi^*(\mu)\,D^*(\mu|\mathbb{\mathrm b}).
\]
Using~\eqref{eq:fisher_0^1}, we write
\[
\Psi(\mu)
=\iiint \frac{1}{t x^2+(1-t)y^2}\,dt\,d\mu(x)d\mu(y).
\]
Hence
\[
|\Psi(\mu)-1|
=\left|
\iiint 
\frac{1-(t x^2+(1-t)y^2)}{t x^2+(1-t)y^2}
\,dt\,d\mu(x)d\mu(y)
\right|.
\]

Applying Cauchy--Schwarz with respect to the product measure
$dt\,d\mu(x)d\mu(y)$ gives
\[
|\Psi(\mu)-1|
\le
\left(
\iiint \frac{1}{(t x^2+(1-t)y^2)^2}
\right)^{1/2}
\left(
\iiint (1-(t x^2+(1-t)y^2))^2
\right)^{1/2}.
\]

For the second factor,
\[
\iiint(1-W_t)^2\,dt\,d\mu\,d\mu = \int_0^1\mathrm{Var}(W_t)\,dt,
\]
where $W_t=tX^2+(1-t)Y^2$. Since $\mathrm{Var}(W_t)=[t^2+(1-t)^2]D^*(\mu|\mathrm{b})^2$ and $t^2+(1-t)^2\leq 1$, we obtain
\[
\int_0^1\mathrm{Var}(W_t)\,dt \leq D^*(\mu|\mathrm{b})^2,
\]
so the second factor is at most $D^*(\mu|\mathrm{b})$.

The first factor equals $\Psi^*(\mu)$, because
\[
\iiint \frac{1}{(t x^2+(1-t)y^2)^2}
=
\iint \frac{1}{x^2-y^2}\!\left(\frac{-1}{x^2}-\frac{-1}{y^2}\right)
d\mu(x)d\mu(y)
=\Psi^*(\mu)^2.
\]
Thus
\[
|\Psi(\mu)-1|
\le \Psi^*(\mu)\,D^*(\mu|\mathbb{\mathrm b}).
\]

Since
\[
\Psi(\mu|\mathbb{\mathrm b})
=\Psi(\mu)+m_2(\mu)-2,
\]
we obtain
\[
\Psi(\mu|\mathbb{\mathrm b})
\le \Psi^*(\mu)\,D^*(\mu|\mathbb{\mathrm b})
+|m_2(\mu)-1|.
\]
Applying this to $\mu_n$, we have $m_2(\mu_n)=1$, while
$\Psi^*(\mu_n)=\Psi^*(\mu)$ by Theorem~\ref{th:constancy}
and
$D^*(\mu_n|\mathbb{\mathrm b})=\frac{1}{\sqrt n}\sqrt{m_4(\mu)-1}$
by Proposition~\ref{prop:Berryesseen}.
This yields the claimed bound.
\end{proof}

\end{document}
